% ---
% Pacotes básicos 
% --
\usepackage{times} %{lmodern}	% Usa a fonte Latin Modern			
\usepackage[T1]{fontenc}% Selecao de codigos de fonte.
\usepackage[utf8]{inputenc}% Codificacao do documento (conversão automática dos acentos)
\usepackage{indentfirst}	% Indenta o primeiro parágrafo de cada seção.
\usepackage[top=30mm, bottom=20mm, left=30mm, right=20mm]{geometry}
\usepackage{color}		% Controle das cores
\usepackage{graphicx}	% Inclusão de gráficos
\graphicspath{ {./build/figuras/} }
\usepackage{microtype} 	% para melhorias de justificação
\usepackage{fancyhdr} %define o estilo do documento
\usepackage{array} %permite o uso de tabelas
\usepackage{longtable} %permite o uso de tabelas personalizadas
\usepackage{rotating} %permite rotacionar imagens
\usepackage{tikz} %permite o desenho de formas geométricas simples dadas suas características
\usepackage{PPGAN}
\usepackage{ragged2e}
\usepackage{csquotes}
\usepackage{pdfpages}
\usepackage{amsmath}
\usepackage{amssymb}
% ---
% Pacotes de citações
% ---
\usepackage[style=abnt, repeatfields, backend=biber]{biblatex}
\addbibresource{referencias.bib} % Seus arquivos de

\ifthenelse{\equal{\ABNTEXisarticle}{true}}{%
\renewcommand{\maketitlehookb}{}
}{}
\PassOptionsToPackage{brazilian,english}{babel}

\DeclareUnicodeCharacter{031C}{ç}